\chapter{Flujo de Ingresos Proyectados}

Para ingresar al mercado, una estrategia razonable sería ofrecer el servicio a un precio competitivo dentro de los rangos más accesibles actualmente disponibles. De acuerdo con la Tabla 4.1, se ha decidido establecer un costo de \$0.36 por cada 1000 mensajes procesados.

Actualmente, entre los clientes públicos de este tipo de servicios se encuentran plataformas de medios digitales como \textit{El País} y \textit{The New York Times}, las cuales pueden generar diariamente en promedio alrededor de 200 contenidos periodísticos, acompañados de más de 400 comentarios por publicación.

Para efectos del presente análisis, se considerará un escenario hipotético en el cual las plataformas de medios digitales nacionales recibirán un volumen equivalente al que muestra la plataforma \textit{The New York Times} (tener en cuenta que la plataforma internacional recibe más mensajes de los que se ha mencionado, solo que estos no son publicados porque ya fueron moderados). Además se considera un crecimiento proyectado según la CAGR del 13,4\%, compatible con el mercado de soluciones de moderación. Estos valores se muestran en la Tabla 5.1.

\begin{table}[h!]
    \centering
    \begin{tabular}{l l}
        \hline
        \textbf{Parámetro} & \textbf{Valor} \\
        \hline
        Precio por 1.000 mensajes & \$0.36 USD \\
        Mensajes procesados diarios & 80.000 mensajes \\
        Ingresos diarios & 28.8 USD \\
        Ingresos mensuales & 864 USD \\
        Mensajes procesados anuales & 29,200,000 mensajes \\
        Tasa de crecimiento anual (CAGR) & 13.4\% \\
        \hline
    \end{tabular}
    \caption{Estimaciones por Servicio \textbf{[Elaboración Propia]}}
\end{table}

Bajo estas consideraciones, la Tabla 5.2 y la Figura 5.1 presentan la proyección de ingresos anuales por cliente para el periodo 2026--2031. Se parte del supuesto de que el precio del servicio permanece constante en \$0.36 USD por cada 1,000 mensajes procesados, y que el crecimiento anual en el uso del servicio por parte de un cliente sigue una CAGR del 13.4\%, coherente con el crecimiento estimado del mercado global de soluciones de moderación de contenido.









\vspace{0.5cm}

\begin{table}[h!]
    \centering
    \begin{tabular}{l c c}
        \hline
        \textbf{Año} & \textbf{Miles de mensajes aprox.} & \textbf{Ingresos} \\
        \hline
        2026 & 29,200 & \$10,512.00 \\
        2027 & 33,113 & \$11,920.61 \\
        2028 & 37,550 & \$13,517.97 \\
        2029 & 42,582 & \$15,329.38 \\
        2030 & 48,288 & \$17,383.51 \\
        2031 & 54,758 & \$19,712.90 \\
        \hline
    \end{tabular}
    \caption{Proyección de ingresos para un cliente \textbf{[Elaboración Propia]}}
\end{table}

\begin{figure}[h!]
    \centering
    % \includegraphics[width=0.72\textwidth]{figura_5_1.png}
    \caption{Flujo de Ingresos para un cliente \textbf{[Elaboración Propia]}}
\end{figure}

Bajo este escenario, se proyecta que un cliente que inicialmente procesa 29,200 millones de mensajes en 2026 aumentará progresivamente su volumen de uso, alcanzando aproximadamente 54.76 millones de mensajes en 2031. Como consecuencia, el ingreso anual generado por dicho cliente pasaría de \$10,512 USD en 2026 a \$19,712 USD en 2031.

Esta proyección refleja un crecimiento sostenido en el uso del servicio, lo cual podría derivarse del aumento en la actividad digital de los clientes, una mayor necesidad de moderación automatizada, o la consolidación del servicio como parte integral de sus operaciones.

En la Figura 5.2 se asume que el número de clientes también aumenta y además, el volumen de mensajes procesados por cada cliente aumenta anualmente siguiendo la CAGR del 13.4\%, alineada con el crecimiento estimado del mercado de soluciones de moderación de contenido.








\vspace{0.5cm}

\begin{figure}[h!]
    \centering
    % \includegraphics[width=0.75\textwidth]{figura_5_2.png}
    \caption{Flujo de Ingresos Totales \textbf{[Elaboración Propia]}}
\end{figure}

Este enfoque refleja una evolución más realista del negocio en una etapa de expansión:

\begin{itemize}
    \item Por un lado, el número de clientes crece conforme más empresas adoptan servicios automatizados de moderación de lenguaje.
    \item Por otro lado, cada cliente tiende a procesar un volumen cada vez mayor de mensajes, a medida que integran más áreas de su plataforma o crecen en usuarios.
\end{itemize}






